\section*{Chapter 13 --- Bring-up}
\addcontentsline{toc}{section}{Chapter 13 --- Bring-up}

\solution{ex:13:1}{The ladder}
\ans{a} Rung 1 proves that \emph{your} SPI master works: the clock, the bit order, the timing and
the configuration. The FPGA is not connected and cannot affect the outcome. Rung 2 proves that the
whole chain works --- the wiring, the slave's synchronisation, the transaction bridge and the
register bank --- and that both sides interpret the command byte the same way.
\ans{b} Because a loopback also gives back what you wrote. If you write \code{0x00000123} and read
\code{0x00000123}, that may just as well be your own \code{MOSI} coming back on \code{MISO}
through a miswiring. \textbf{The masking is the proof:} write \code{0xFFFFFFFF} to \code{TX_ID}
and read back \code{0x000007FF}. An echo cannot mask away the upper bits; only a register bank
can.
\ans{c} Rung 1. \code{AvrSpiTransport} is the only class in the project that no automated test
covers, since it talks to a physical register --- it cannot be run on the host computer, and a
test double of it would have tested the double.

\solution{ex:13:2}{From symptom to hypothesis}
\ans{a} For example: \code{MISO} not connected or connected to the wrong pin; the pin mux not set,
so the peripheral sits on the wrong port; \code{VDDIO2} not powered, so \code{PORTC} does not
work; the FPGA not programmed; the clock too fast; \code{SS} not released between transactions; or
a missing common ground.
\ans{b} Fastest first: \code{VDDIO2} (read the MVIO status in the program --- no measurement at
all), the pin mux (read the code), the FPGA being programmed (look at the board), then the
measurements: clock pulses on \code{SCK}, the \code{SS} behaviour, and last the wiring, wire by
wire.
\ans{c} A loopback: connect \code{MOSI} straight to \code{MISO} and run rung 1 again. If you get
back what you sent, your master works and the fault is in the wiring to the FPGA or in the FPGA.
If you do not, the fault is on your side, and the FPGA is ruled out.

\solution{ex:13:3}{Whose fault?}
\ans{a} It cannot be decided from the symptom alone --- all three are possible, which is the whole
difficulty of debugging a system two groups built one half each of.
\ans{b} Climb down the ladder until something gives the right receipt. If rungs 2 and 3 give the
right answers the SPI chain works, so the fault is on the CAN side: the wiring to the transceiver,
the termination, or the controller. If rung 3 gives \code{0x00000001} but \code{STATUS} bit 0
never goes low after a \code{TX_SEND} write, the trigger is not reaching the controller --- which
points at the register bank, or at your write not having bit 0 set.
\ans{c} It should contain what you did, what you observed and what you have already ruled out. For
example: \emph{"Rung 2 gives the right masking for \code{TX_ID} (\code{0xFFFFFFFF} reads back as
\code{0x000007FF}), so the SPI chain and the register bank work. \code{STATUS} reads as
\code{0x00000001}. After a write of \code{0x1} to \code{TX_SEND}, bit 0 never goes low, and
nothing shows on \code{tx_bus}. We have verified that our command byte is \code{0x85} and that bit
0 is set in the data."} That is a report somebody can act on; "it does not work" is not.

\solution{ex:13:4}{The termination}
\ans{a} About 60~\textohm: two 120-ohm resistors in parallel, one at each end.
\ans{b} There are three resistors instead of two --- three 120~\textohm{} in parallel give
40~\textohm. Most likely a node in the middle of the bus has its termination switched on.
\ans{c} There is only one resistor; the other is missing or not connected. It shows when the cable
gets longer, the speed higher or the environment noisier --- that is, typically when the design is
moved off the lab bench, which is the worst possible moment.
